%% kapa: exam=1 round=33 year=2022 period=1 form=A subject=civil_law subject_ko=민법 seq=1 number=1
%% source: Q-Net 감정평가사 기출문제 (cst00309) artlSeq=5209997 file=raw/1008/5209997/2022년도 제33회 감정평가사 1차시험 1교시(민법,경제학,부동산학).pdf
%% conversion: kordoc 4.24.51 -> tools/gen1.py (auto); review=proofread-sonnet; flags=none
\begin{question}{1}
민법의 법원(法源)에 관한 설명으로 옳은 것을 모두 고른 것은? (다툼이 있으면 판례에 따름)
\begin{framedtext}
ㄱ. 헌법에 의해 체결·공포된 민사에 관한 조약은 민법의 법원이 되지 않는다. \par
ㄴ. 관습법이 되기 위해서는 사회구성원의 법적 확신이 필요하다. \par
ㄷ. 관습법은 법령에 저촉되지 않는 한 법칙으로서의 효력이 있다.
\end{framedtext}
\begin{choices}
  \item ㄱ
  \item ㄴ
  \item ㄱ, ㄷ
  \item ㄴ, ㄷ
  \item ㄱ, ㄴ, ㄷ
\end{choices}
\end{question}
\answer{4}
